\clearpage\begingroup\NoSourceHead\SourceLeaf{509}
\renewcommand{\thefootnote}{\arabic{footnote})}\setcounter{footnote}{0}
\ArithmeticHead{52}{L}{BOLZA, GAUSS UND DIE VARIATIONSRECHNUNG.}
\begin{center}\bfseries II. Teil: Die Disquisitiones generales circa superficies curvas.\end{center}
\noindent\hspace*{\parindent}1. Die Beziehungen der \textit{Disquisitiones generales circa superficies curvas} zur
Variationsrechnung werden sämtlich durch die Theorie der geodätischen Linien
vermittelt. Es kommen dabei in Betracht
\begin{list}{}{\setlength{\leftmargin}{2.8em}\setlength{\labelwidth}{1.8em}\setlength{\labelsep}{.5em}\setlength{\topsep}{0pt}\setlength{\itemsep}{0pt}\setlength{\parsep}{0pt}}
\item[1)] die Ableitung der Differentialgleichung der geodätischen Linien (Art. 14, 18),
\item[2)] die Sätze über geodätische Polar- und Parallelkoordinaten (Art. 15, 16, 19, 22),
\item[3)] der Satz über die Totalkrümmung eines geodätischen Dreiecks (Art. 20).
\end{list}
\begin{center}\textls[100]{A. Die Differentialgleichung der geodätischen Linien.}\end{center}
\noindent\hspace*{\parindent}2. Die Geschichte der geodätischen Linien vor \textsc{Gauss} ist von \textsc{Stäckel}\footnote{P. \textsc{Stäckel}, \textit{Bemerkungen zur Geschichte der geodätischen Linien}, Leipziger Berichte, Bd. 45 (1893),
S. 444; vgl. auch G. \textsc{Eneström}, \textit{Sur la découverte de l’équation générale des lignes géodesiques}, Bibliotheca
mathematica (2), Bd. 13 (1899), S. 19 und P. \textsc{Stäckel}, \textit{Gauss als Geometer}, Materialien für eine wissenschaftl.
Biographie von Gauss, Heft 5, 1918 und Werke X 2, Abh. IV, Nr. 30; am Schluss der letztgenannten Arbeit
befindet sich auch eine Bibliographie der \textsc{Gauss}schen \textit{Disquisitiones}.},
ausführlich dargestellt worden; wir können uns daher darauf beschränken, diejenigen
Stellen hervorzuheben, an denen \textsc{Gauss} über seine Vorgänger hinausgegangen
ist.

\textsc{Gauss} hat die Differentialgleichung der geodätischen Linien auf zwei verschiedene
Arten abgeleitet, indem er die Fläche das eine Mal durch eine
Gleichung zwischen den Koordinaten gegeben sein lässt (Art. 14), das andere
Mal in Parameterdarstellung annimmt (Art. 19).
\begin{center}\textls[100]{a) Die Fläche ist durch eine Gleichung gegeben.}\end{center}
\noindent\hspace*{\parindent}3. Hier schliesst sich \textsc{Gauss} an \textsc{Lagrange} an. \textsc{Lagrange} hatte in den
\textit{Leçons sur le calcul des fonctions} (Ausgabe von 1806)\footnote{Oeuvres de \textsc{Lagrange}, Bd. X, S. 435.} die Differentialgleichung
der geodätischen Linien aus dem Variationsproblem
\[
\int\sqrt{1+\left(\frac{dy}{dx}\right)^2+\left(\frac{dz}{dx}\right)^2}\,dx=\text{Minimum}
\]
\clearpage\SourceLeaf{510}
\ArithmeticHead{53}{R}{DIE DIFFERENTIALGLEICHUNG DER GEODÄTISCHEN LINIEN.}
\noindent mit der Nebenbedingung
\begin{equation*}\tag{1}
F(x,y,z)=0
\end{equation*}
in der Form erhalten
\[
R\frac{d}{dx}\left(\frac{dy}{ds}\right)-Q\frac{d}{dx}\left(\frac{dz}{ds}\right)=0,
\]
wo
\[
P=\frac{\partial F}{\partial x},\quad Q=\frac{\partial F}{\partial y},\quad R=\frac{\partial F}{\partial z},
\]
und daraus dann die Fundamentaleigenschaft der geodätischen Linien abgeleitet,
dass die Schmiegungsebene der Kurve auf der Tangentialebene der
Fläche senkrecht steht.

\textsc{Gauss} weicht hiervon nur darin ab, dass er die zulässigen Kurven in
Parameterdarstellung annimmt und daher das Längenintegral in der Form
ansetzt
\[
\int\sqrt{\left(\frac{dx}{dw}\right)^2+\left(\frac{dy}{dw}\right)^2+\left(\frac{dz}{dw}\right)^2}\,dw.
\]
Das hat den Vorteil, dass er die Differentialgleichung der geodätischen Linien
in der \textls[100]{symmetrischen} Form
\begin{equation*}\tag{2}
\frac{d^2x}{ds^2}:\frac{d^2y}{ds^2}:\frac{d^2z}{ds^2}=P:Q:R
\end{equation*}
erhält, die unmittelbar die charakteristische Eigenschaft der geodätischen Linien
ausdrückt.

Was die Strenge des Variationsschlusses anbetrifft, so geht \textsc{Gauss} nicht
über \textsc{Lagrange} hinaus. Aus dem Verschwinden der ersten Variation in der
Form
\[
\int\left(\delta x\,d\left(\frac{dx}{ds}\right)+\delta y\,d\left(\frac{dy}{ds}\right)+\delta z\,d\left(\frac{dz}{ds}\right)\right)=0
\]
schliesst er mit einem »constat«, das sich offenbar auf \textsc{Lagrange} bezieht,
\[
\delta x\,d\left(\frac{dx}{ds}\right)+\delta y\,d\left(\frac{dy}{ds}\right)+\delta z\,d\left(\frac{dz}{ds}\right)=0.
\]
Diese Gleichung führt dann mit der aus (1) folgenden Gleichung
\[
P\delta x+Q\delta y+R\delta z=0
\]
auf die Proportion (2).
\par\clearpage\endgroup
